\documentclass[aspectratio=169,10pt]{beamer}
\usetheme[numbering=fraction,progressbar=frametitle]{metropolis}
\usepackage{fontspec}
\setsansfont{Latin Modern Sans}
\usepackage{amsmath,amssymb,mathtools,graphicx,booktabs}
\usepackage{polyglossia}
\setdefaultlanguage{english}
\setbeamertemplate{navigation symbols}{}
\setbeamerfont{frametitle}{size=\large}
\setbeamerfont{normal text}{size=\normalsize}
\setlength{\parskip}{.45em}
\hypersetup{pdfauthor={},pdfsubject={Lecture slides},pdfkeywords={},pdftitle={Coordinate geometry: lines, projections and circles}}
\providecommand{\N}{\mathbb{N}}
\providecommand{\Z}{\mathbb{Z}}
\providecommand{\D}{\mathbb{D}}
\providecommand{\Q}{\mathbb{Q}}
\providecommand{\R}{\mathbb{R}}
\providecommand{\C}{\mathbb{C}}
\providecommand{\vect}[1]{\overrightarrow{#1}}
\providecommand{\vi}{\vec{\imath}}
\providecommand{\vj}{\vec{\jmath}}
\providecommand{\vk}{\vec{k}}
\providecommand{\norm}[1]{\left\lVert #1\right\rVert}
\providecommand{\vc}[2]{\begin{pmatrix}#1\\#2\end{pmatrix}}
\providecommand{\deter}[2]{\det\left(#1,#2\right)}
\providecommand{\abs}[1]{\left\lvert #1\right\rvert}
\providecommand{\dd}{\,\mathrm{d}}
\providecommand{\e}{\mathrm{e}}
\providecommand{\Card}{\operatorname{Card}}
\providecommand{\E}{\operatorname{E}}
\providecommand{\V}{\operatorname{V}}
\providecommand{\degre}{^{\circ}}
\providecommand{\un}[1]{\,\mathrm{#1}}
\providecommand{\sci}[2]{#1\times10^{#2}}
\providecommand{\pgcd}{\operatorname{PGCD}}
\providecommand{\ppcm}{\operatorname{PPCM}}
\title{Coordinate geometry: lines, projections and circles}
\subtitle{Lecture and worked examples}
\date{}
\author{}
\begin{document}
\begin{frame}[plain]
\vfill
{\usebeamerfont{title}\inserttitle\par}
\medskip
{\usebeamerfont{subtitle}\insertsubtitle\par}
\medskip
{\usebeamercolor[fg]{progress bar}\rule{\textwidth}{.5pt}}
\vfill
\end{frame}
\begin{frame}{Chapter path}
\textbf{Objectives.} Use a normal vector; calculate a projection; determine or recognise a circle; translate a configuration into equations.

\textbf{Prerequisites.} Vector coordinates, scalar product, linear systems, completing the square. The frame is always orthonormal.

\textbf{Yearly reference point.} 5, 7 and 12 January: 6 hours
\end{frame}
\begin{frame}{Normal vector and Cartesian equation}
A nonzero normal vector to a line is orthogonal to every direction vector of that line. If \(\vec n=(a,b)\ne(0,0)\) is normal and \(A(x_A,y_A)\) belongs to d, then \(d\) has equation \(a(x-x_A)+b(y-y_A)=0\), that is \(ax+by+c=0\). A possible direction vector is \((-b,a)\).
\end{frame}
\begin{frame}{Normal vector and Cartesian equation}
Through \(A(1,2)\) with normal \((2,-3)\): \(2(x-1)-3(y-2)=0\), so \(2x-3y+4=0\). The vector \((3,2)\) directs the line: product \(6-6=0\).
\end{frame}
\begin{frame}{Normal vector and Cartesian equation}
\((a,b)\) is normal, not a direction vector in general. \((a,b)\) must not be zero.
\end{frame}
\begin{frame}{Normal vector and Cartesian equation}
Equation of the line through \((2,-1)\) with normal \((3,4)\).
\end{frame}
\begin{frame}{Orthogonal projection and distance}
The projection H of P onto d is characterised by \(H\in d\) and \(\vect{PH}\) collinear with a normal vector of d. For \(d:ax+by+c=0\), write \(H=P-\lambda(a,b)\), with \(\lambda =\dfrac{ax_P+by_P+c}{a^2+b^2}\). The distance is \(PH=\dfrac{|ax_P+by_P+c|}{\sqrt {a^2+b^2}}\). The distance formula is a deduced supplement, not an isolated routine.
\end{frame}
\begin{frame}{Orthogonal projection and distance}
\begin{columns}[T,onlytextwidth]
\begin{column}{.57\textwidth}
For \(P(3,4)\) and \(d:x+y-2=0\), \(\lambda =\dfrac{3+4-2}{2}=\dfrac{5}{2}\), so \(H(\dfrac{1}{2},\dfrac{3}{2})\) and \(PH=\dfrac{5}{\sqrt 2}=\dfrac{5\sqrt 2}{2}\). Check \(x_H+y_H=2\).
\end{column}
\begin{column}{.4\textwidth}
\centering
\includegraphics[width=\linewidth,height=4.5cm,keepaspectratio]{figures/premiere-07-1-en.pdf}
\end{column}
\end{columns}
\end{frame}
\begin{frame}{Orthogonal projection and distance}
The distance is positive or zero: keep the absolute value. Two checks, membership and perpendicularity, secure the projection.
\end{frame}
\begin{frame}{Orthogonal projection and distance}
Project \(P(4,1)\) onto \(d:x=2\).
\end{frame}
\begin{frame}{Equation of a circle}
The circle with centre \(I(\alpha,\beta)\) and radius \(r>0\) has equation \((x-\alpha)^2+(y-\beta)^2=r^2\). An equation \(x^2+y^2+ux+vy+w=0\) is handled by completing both squares. If the squared radius obtained is negative, the set is empty; if it is zero, it is a point.
\end{frame}
\begin{frame}{Equation of a circle}
\(x^2+y^2-4x+6y+9=0\) becomes \((x-2)^2+(y+3)^2=4\). Centre \((2,-3)\), radius 2. The signs of the coordinates of the centre are opposite to those appearing in the brackets.
\end{frame}
\begin{frame}{Equation of a circle}
A quadratic equation in x and y is not always a circle: check the coefficients and the real radius.
\end{frame}
\begin{frame}{Equation of a circle}
Recognise \(x^2+y^2+2x-8y+17=0\).
\end{frame}
\begin{frame}{Solving a configuration}
An intersection is obtained by solving the equations simultaneously. For a circle and a line parallel to an axis, substituting the constant coordinate gives a quadratic equation. For a perpendicular bisector, write the equality of the squared distances; the quadratic terms cancel.
\end{frame}
\begin{frame}{Solving a configuration}
The circle \(x^2+y^2=25\) and the line \(x=3\) meet at \((3,4)\) and \((3,-4)\). The perpendicular bisector of \(A(0,0)\) and \(B(4,2)\) satisfies \(x^2+y^2=(x-4)^2+(y-2)^2\), so \(2x+y=5\).
\end{frame}
\begin{frame}{Solving a configuration}
An intersection is not the union of the solution sets. Check the points in both equations.
\end{frame}
\begin{frame}{Solving a configuration}
Intersect \((x-1)^2+(y+2)^2=9\) with \(y=1\).
\end{frame}
\end{document}
